\chapter{Meccanica quantistica relativistica}
Ora analizzeremo come la meccanica quantistica può essere adattata per descrivere fenomeni in regime relativistico. La nota equazione di Schrödinger:
\begin{equation}
    i \hbar \pdv{\psi}{t}  = \hat{H} \psi
\end{equation}
Può descrivere il comportamento di numerosi sistemi subatomici, tuttavia non è valida quando si entra in regimi di velocità relativistiche. Dunque serve un'equazione che sia consistente con il formalismo della meccanica quantistica e rispetti le leggi della relatività speciale.
\section{Equazione di Klein-Gordon}

\addfigure[Oskar Klein. Fonte: \href{https://en.wikipedia.org/wiki/Oskar_Klein}{Wikipedia}]{jpg/Oskar_Klein}{0.3}

Per raggiungere questo obiettivo possiamo ipotizzare di applicare la regola di quantizzazione canonica e trovare i seguenti operatori differenziali:
\begin{equation}
    \begin{split}
        &\vec{p} \rightarrow -i\grad \\[8pt]
        &E \rightarrow i\pdv{}{t}
    \end{split}
\end{equation}
\textbf{N.B.} Da questo momento utilizzeremo le unità di misura naturali ovvero poniamo $c = \hbar = 1$, da qui deriva l'assenza di altri termini negli operatori.
Una volta individuati i ``sostituti'' di energia e quantità di moto possiamo inserirli nell'equazione relativistica per l'energia:
\begin{equation}
    \begin{split}
        &E^2 = p^2 + m^2 \\[8pt]
        &-\pdv[2]{}{t} = -\lap + m^2 \\[8pt]
        &\pdv[2]{}{t} -\lap + m^2 = 0
    \end{split}
\end{equation}
Definendo l'operatore d'onda o D'Alembertiano come:
\begin{equation}
    \dalop = \pdv[2]{}{t} -\lap
\end{equation}
Otteniamo un operatore che possiamo applicare ad una funzione d'onda ottenendo l'equazione di \textbf{Klein-Gordon} (KG).
\begin{theorem}{Equazione di Klein-Gordon}
    L'equazione relativistica per una particella quantistica libera di massa $m$ è data da:
    \begin{equation}
        \brackets{\dalop + m^2}\psi = 0
    \end{equation}
\end{theorem}
O, alternativamente in notazione tensoriale:
\begin{equation}
    \boxed{\brackets{\partial^\mu\partial_\mu + m^2} \psi = 0}
\end{equation}
Ricordando che per la derivazione sono state utilizzate le unità naturali, possiamo riscrivere l'equazione di KG nelle unità S.I. come:
\begin{equation}
    \boxed{\brackets{\dfrac{1}{c^2}\pdv[2]{}{t} - \lap + \dfrac{m^2 c^2}{\hbar^2}} \psi = 0}
\end{equation}
Poiché abbiamo costruito questa equazione mediante la formula per l'energia relativistica, che si può ottenere come prodotto scalare del tensore $p^\mu$, otteniamo che l'equazione di Klein-Gordon è invariante per trasformazioni di Lorentz come desiderato.\\
Ipotizziamo ora una soluzione per l'equazione del tipo onda piana:
\begin{equation}
    \psi = N\efunction^{-i\brackets{\vec{p}\cdot\vec{r} - Et}}
\end{equation}
Sostituendo all'interno dell'equazione di KG otteniamo:
\begin{equation}
    \begin{split}
        \lap \psi &= -\vec{p}^2\psi \\[8pt]
        \pdv[2]{\psi}{t}  &= E^2\psi \\[8pt]
        &\Downarrow \\[8pt]
        E^2 &= p^2 + m^2
    \end{split}
\end{equation}
\section{Equazione di Dirac}

\addfigure[Paul Dirac. Fonte: \href{https://en.wikipedia.org/wiki/Paul_Dirac}{Wikipedia}]{jpg/Paul_Dirac}{0.3}

L'equazione di KG è ancora oggi fondamentale, tuttavia è valida solo per particelle a spin nullo, infatti, già nel formalismo di Schrödinger lo spin viene inserito come contributo ``anomalo'' all'interno della spiegazione di molti fenomeni quantistici.
Paul Dirac dunque vuole ottenere un'equazione che rispetti le leggi relativistiche, e dunque sia consistente con l'equazione di KG, ma che contenga anche informazioni valide per descrivere particelle fondamentali come gli elettroni.
La brillante idea che viene utilizzata per ricavare la famosa equazione consiste nel ricercare una relazione lineare tra gli operatori differenziali presenti. In particolare abbiamo:
\begin{equation}
    \begin{split}
        \hat{H}\psi &= \brackets{\vec{\alpha} \cdot \vec{p} + \beta m}\psi \\[8pt]
        &= \brackets{-i\alpha_x\pdv{}{x} -i\alpha_y\pdv{}{y} -i\alpha_z\pdv{}{z} + \beta m}\psi = -i\pdv{\psi}{t}
    \end{split}
\end{equation}
Ipotizziamo di applicare l'operatore a se stesso:
\begin{equation}
    \begin{split}
        &\left[ -i\alpha_x\pdv{}{x} -i\alpha_y\pdv{}{y} -i\alpha_z\pdv{}{z} + \beta m \right]^2  = \\[8pt]
        &= \left[ -i\alpha_x\pdv{}{x} -i\alpha_y\pdv{}{y} -i\alpha_z\pdv{}{z} + \beta m \right]
            \left[ -i\alpha_x\pdv{}{x} -i\alpha_y\pdv{}{y} -i\alpha_z\pdv{}{z} + \beta m \right]  \\[8pt]
        &=
                \brackets{-i\alpha_x\pdv{}{x}}^2 + \brackets{-i\alpha_y\pdv{}{x}}^2 + \brackets{-i\alpha_z\pdv{}{x}}^2
                + (\beta m)^2
                + \text{cross terms}
        \\[8pt]
        &=
                -\alpha_x^2 \pdv[2]{}{x}
                -\alpha_y^2 \pdv[2]{}{y}
                -\alpha_z^2 \pdv[2]{}{z}
                + \beta^2 m^2
                + \text{cross terms}
        \\[8pt]
    \end{split}
\end{equation}
I termini incrociati sono:
\begin{equation}
    \begin{split}
        &\left[
                \brackets{-i\alpha_x\pdv{}{x}} \brackets{-i\alpha_y\pdv{}{y}}
                + \brackets{-i\alpha_x\pdv{}{x}} \brackets{-i\alpha_z\pdv{}{z}}
                + \brackets{-i\alpha_y\pdv{}{y}} \brackets{-i\alpha_z\pdv{}{z}}
                + \text{simmetrici}
        \right] \\[8pt]
        &= -\left[
                \alpha_x\alpha_y \pdv{}{x}{y}
                + \alpha_x\alpha_z \pdv{}{x}{z}
                + \alpha_y\alpha_z \pdv{}{y}{z}
                + \text{simmetrici}
        \right] \\[8pt]
    \end{split}
\end{equation}
Inoltre, i termini misti con $\beta m $ sono:
\begin{equation}
    -im \sqbr{\brackets{\beta \alpha_x + \alpha_x \beta}\pdv{}{x} + \brackets{\beta \alpha_y + \alpha_y \beta}\pdv{}{y} + \brackets{\beta \alpha_z + \alpha_z \beta}\pdv{}{z}}\\[8pt]
\end{equation}
Raccogliendo tutti i termini:
\begin{equation}
    \begin{split}
        &= \bigg[
                -\sum_{j} \alpha_j^2 \pdv[2]{}{x_j}
                - \sum_{j<k} (\alpha_j\alpha_k + \alpha_k\alpha_j) \pdv{}{x_j}{x_k}\\[8pt]
                &+ \beta^2 m^2
                - i m \sum_j \brackets{\beta \alpha_j + \alpha_j\beta} \pdv{}{x_j}
        \bigg]
    \end{split}
\end{equation}
Per ottenere l'equazione di KG dobbiamo imporre le seguenti condizioni:
\begin{equation}
    \begin{cases}
        \alpha_j^2 = 1 \\[8pt]
        \beta^2 = 1 \\[8pt]
        \alpha_j\alpha_k + \alpha_k\alpha_j = 0 \quad (j \neq k) \\[8pt]
        \beta\alpha_j + \alpha_j\beta = 0
    \end{cases}
\end{equation}
Dalle condizioni 3 e 4 osserviamo che $\alpha_j$ e $\beta$ non possono essere semplici numeri, infatti la relazione che devono rispettare è del tipo:
\begin{equation}
    AB = - BA
\end{equation}
Ovvero $\alpha_j$ e $\beta$ sono matrici anti-commutative. Le richieste 1 e 2 si traducono dunque in:
\begin{equation}
    \begin{cases}
        \alpha_j^2 = \mathbb{I} \\[8pt]
        \beta^2 = \mathbb{I} \\[8pt]
    \end{cases}
\end{equation}
Rimane da determinare la dimensione di queste matrici. Per arrivare a questa conclusione abbiamo bisogno di ricavare altre proprietà interessanti di $\alpha_j$ e $\beta$. Sappiamo che il quadrato di queste matrici è l'identità dunque la traccia di una qualsiasi di esse risulta:
\begin{equation}
    \Tr (\alpha_j) = \Tr (\alpha_j \mathbb{I}) = \Tr (\alpha_j \beta^2) = \Tr (\beta \alpha_j \beta) = -\Tr (\alpha_j \beta^2) = - \Tr (\alpha_j)
\end{equation}
Abbiamo dunque sfruttato l'invarianza della traccia per permutazioni cicliche e la proprietà 4 che richiediamo, ottenendo che la traccia sia nulla $\Tr (\alpha_j) = 0$.
Inoltre, sempre dall'equivalenza con l'identità sappiamo che se $\lambda$ è autovalore di $\alpha_j$ allora $\lambda^2$ è autovalore del quadrato di $\alpha_j$ ovvero dell'identità, dunque:
\begin{equation}
    \lambda^2 = 1 \implies \lambda = \pm 1
\end{equation}
Combinando questa informazione con il fatto che la traccia deve essere nulla escludiamo tutte le matrici di dimensioni dispari, consistente con l'impossibilità di soddisfare l'equazione con numeri (matrici $1 \by 1$).
Le matrici di Pauli rispettano tutte le proprietà sopra citate, tuttavia necessitiamo di 4 matrici, dunque la prima dimensione disponibile è $4 \by 4$. Nonostante le matrici di Pauli non siano direttamente le matrici da noi richieste possiamo sfruttare le loro proprietà per costruire facilmente $\alpha_j$ e $\beta$.
Osserviamo inoltre che, così facendo la funzione d'onda $\psi$ deve essere composta di almeno $4$ elementi:
\begin{equation}
    \psi = \begin{pmatrix}
    \psi_1 \\
    \psi_2 \\
    \psi_3 \\
    \psi_4
    \end{pmatrix}
\end{equation}
Queste particolari funzioni d'onda sono dette \textbf{spinori}.\\
Come precedentemente detto ci servono 4 matrici hermitiane e legate da una relazione anti-commutativa. Possiamo dunque costruire queste matrici mediante il prodotto tensoriale tra le matrici di Pauli. In particolare abbiamo che:
\begin{equation}
  \alpha_i = \begin{pmatrix}
    0 & \sigma_i \\
    \sigma_i & 0
  \end{pmatrix}, \qquad
  \beta = \begin{pmatrix}
    \mathbb{I} & 0 \\
    0 & -\mathbb{I}
  \end{pmatrix}
\end{equation}
Osserviamo immediatamente che queste matrici rispettano le relazioni desiderate grazie alle proprietà delle matrici di Pauli, infatti:
\begin{equation}
  \alpha_i \alpha_j =
  \begin{pmatrix}
    0 & \sigma_i \\
    \sigma_i & 0
  \end{pmatrix}
  \begin{pmatrix}
    0 & \sigma_j \\
    \sigma_j & 0
  \end{pmatrix} =
  \begin{pmatrix}
    \sigma_i\sigma_j & 0 \\
    0 & \sigma_i\sigma_j
  \end{pmatrix} =
  \begin{pmatrix}
    -\sigma_j\sigma_i & 0 \\
    0 & -\sigma_j\sigma_i
  \end{pmatrix} = -\alpha_j \alpha_i
\end{equation}
Ora analizziamo cosa capita alla corrente di densità di probabilità. Abbiamo che:
\begin{equation}
  -i\alpha_x\pdv{\psi}{x} -i\alpha_y\pdv{\psi}{y} -i\alpha_z\pdv{\psi}{z} + m\beta\psi = i\pdv{\psi}{t}
\end{equation}
E la sua complessa coniugata, ricordando l'hermitianità delle matrici $\alpha$ e $\beta$ risulta:
\begin{equation}
  i\pdv{\psi^\dagger}{x}\alpha_x + i\pdv{\psi^\dagger}{y}\alpha_y + i\pdv{\psi^\dagger}{z}\alpha_z + m\psi^\dagger\beta = -i\pdv{\psi^\dagger }{t}
\end{equation}
Da cui:
\begin{equation}
  \begin{split}
    &\psi^\dagger\brackets{-i\alpha_x\pdv{\psi}{x} -i\alpha_y\pdv{\psi}{y} -i\alpha_z\pdv{\psi}{z} + m\beta\psi} - \brackets{i\pdv{\psi^\dagger}{x}\alpha_x + i\pdv{\psi^\dagger}{y}\alpha_y + i\pdv{\psi^\dagger}{z}\alpha_z + m\psi^\dagger\beta}\psi = \\[8pt]&= i\psi^\dagger\pdv{\psi}{t}-i\pdv{\psi^\dagger }{t}\psi
  \end{split}
\end{equation}
Ora, osservando che:
\begin{equation}
  \psi^\dagger\alpha_x\pdv{\psi}{x} + \pdv{\psi^\dagger}{x}\alpha_x\psi= \pdv{\brackets{\psi^\dagger\alpha_x\psi}}{x}
\end{equation}
e che:
\begin{equation}
  \psi^\dagger\pdv{\psi}{t} + \pdv{\psi^\dagger}{t}\psi = \pdv{\brackets{\psi^\dagger\psi}}{t}
\end{equation}
Possiamo scrivere in forma compatta:
\begin{equation}
  \divergence{\brackets{\psi^\dagger\vecsymb{\alpha}\psi}} = \pdv{\rho}{t}
\end{equation}
E identifichiamo dunque la corrente di densità di probabilità con:
\begin{equation}
  \vec{j} = \psi^\dagger\vecsymb{\alpha}\psi
\end{equation}
Inoltre la densità di probabilità è data da:
\begin{equation}
  \rho = \psi^\dagger \psi = \begin{pmatrix}
    \psi_1^* & \psi_2^* & \psi_3^* & \psi_4^*
  \end{pmatrix}
  \begin{pmatrix}
    \psi_1 \\ \psi_2 \\ \psi_3 \\ \psi_4
  \end{pmatrix} = \abs{\psi_1}^2 + \abs{\psi_2}^2 + \abs{\psi_3}^2 + \abs{\psi_4}^2 > 0
\end{equation}
L'equazione di Dirac dunque riporta sempre una densità di probabilità positiva.
Possiamo ora cercare di riscrivere l'equazione di Dirac in notazione tensoriale. Definiamo le matrici gamma come:
\begin{equation}
  \gamma^0 = \beta, \qquad \gamma^i = \beta \alpha_i
\end{equation}
In questo modo, moltiplicando ambo i membri dell'equazione di Dirac per $\beta$ otteniamo la forma più nota dell'equazione.
\begin{theorem}{Equazione di Dirac}
    L'equazione relativistica per una particella quantistica libera di massa $m$ e spin $\dfrac{1}{2}$ è data da:
    \begin{equation}
        \boxed{\brackets{i\gamma^\mu \partial_\mu - m}\psi = 0}
    \end{equation}
\end{theorem}
Possiamo osservare diverse proprietà delle matrici gamma. In particolare:
\begin{equation}
  \begin{split}
    &\brackets{\gamma^0}^2 = \mathbb{I}\\[8pt]
    &\brackets{\gamma^i}^2 = -\mathbb{I}
  \end{split}
\end{equation}
Da cui otteniamo la seguente relazione di anticommutazione:
\begin{equation}
  \{\gamma^\mu, \gamma^\nu\} = \gamma^\mu \gamma^\nu + \gamma^\nu \gamma^\mu = 2 g^{\mu \nu}
\end{equation}
Questo serve a definire l'algebra che caratterizza le matrici gamma.\\
Inoltre, per definizione le matrici $\alpha_i$ e $\beta$ sono anticommutative e hermitiane, dunque le matrici gamma avranno la seguente proprietà:
\begin{equation}
  \brackets{\gamma^0}^\dagger = \gamma^0, \qquad \brackets{\gamma^i}^\dagger = \brackets{\beta\alpha_i}^\dagger = \alpha_i^\dagger\beta^\dagger = \alpha_i\beta = -\beta \alpha_i = -\gamma^i
\end{equation}
Ovvero la matrice $\gamma^0$ è hermitiana, mentre le altre matrici gamma sono anti-hermitiane.\\
Grazie ai risultati precedenti possiamo anche definire un quadrivettore di corrente come:
\begin{equation}
  j^\mu = \brackets{\rho, \vec{j}} = \psi^\dagger \gamma^0\gamma^\mu \psi
\end{equation}
In questo modo l'equazione di continuità può essere riscritta come:
\begin{equation}
  \partial_\mu j^\mu = 0
\end{equation}
Possiamo ulteriormente semplificare la notazione definendo lo spinore aggiunto come:
\begin{equation}
  \bar{\psi} = \psi^\dagger \gamma^0 = \begin{pmatrix}
    \psi_1^* & \psi_2^* & \psi_3^* & \psi_4^*
  \end{pmatrix} \begin{pmatrix}
    1 & 0 & 0 & 0 \\
    0 & 1 & 0 & 0 \\
    0 & 0 & -1 & 0 \\
    0 & 0 & 0 & -1
  \end{pmatrix} = \begin{pmatrix}
    \psi_1^* & \psi_2^* & -\psi_3^* & -\psi_4^*
  \end{pmatrix}
\end{equation}
Dunque la corrente di densità di probabilità può essere riscritta come:
\begin{equation}
  j^\mu = \bar{\psi} \gamma^\mu \psi
\end{equation}
Ora che abbiamo trovato una nuova equazione invariante per trasformazioni di Lorentz possiamo chiederci se il momento angolare rappresenti ancora un operatore conservato. Il momento angolare è dato da:
\begin{equation}
  \vec{L} = \vec{r} \cross \vec{p}
\end{equation}
Analizziamo per componenti:
\begin{equation}
  L_i = \epsilon_{ijk} x_j p_k
\end{equation}
Calcoliamo il commutatore tra l'hamiltoniana di Dirac e la componente $L_i$:
\begin{equation}
  \begin{split}
    [H, L_i] &= \sqbr{\alpha_j p_j + \beta m, \epsilon_{ikl} x_k p_l } \\[8pt]
    &= \alpha_j \epsilon_{ikl} \sqbr{p_j, x_k p_l} + \cancel{\sqbr{\beta m, \epsilon_{ikl} x_k p_l}} \\[8pt]
    &= \alpha_j \epsilon_{ikl} \brackets{\sqbr{p_j, x_k} p_l +\cancel{ x_k \sqbr{p_j, p_l}}} \\[8pt]
    &= -i \alpha_j \epsilon_{ikl} \delta_{jk} p_l \\[8pt]
    &= -i \epsilon_{ijl} \alpha_j p_l \neq 0
  \end{split}
\end{equation}
Abbiamo quindi ottenuto che:
\begin{equation}
  [H, \vec{L}] = -i\vecsymb{\alpha} \cross \vec{p} \neq 0
\end{equation}
Ovvero che l'hamiltoniana di Dirac non commuta con il momento angolare orbitale. Dunque il momento angolare orbitale non è una quantità conservata in questo formalismo.\\
Identifichiamo ora un nuovo operatore di spin:
\begin{equation}
  \vec{S} = \dfrac{1}{2} \vecsymb{\Sigma} = \dfrac{1}{2} \begin{pmatrix}
    \vecsymb{\sigma} & 0 \\
    0 & \vecsymb{\sigma}
  \end{pmatrix}
\end{equation}
Calcoliamo ora il commutatore tra l'hamiltoniana di Dirac e lo spin:
\begin{equation}
  \begin{split}
    [H, S_i] &= \sqbr{\alpha_j p_j + \beta m, \dfrac{1}{2} \Sigma_i} \\[8pt]
    &= \dfrac{1}{2} \sqbr{\alpha_j p_j, \Sigma_i} + \cancel{\dfrac{1}{2}\sqbr{\beta m, \Sigma_i}} \\[8pt]
    &= \cancel{\dfrac{1}{2} \alpha_j \sqbr{p_j, \Sigma_i}} + \dfrac{1}{2} \sqbr{\alpha_j, \Sigma_i} p_j \\[8pt]
    &= \dfrac{1}{2} \sqbr{\alpha_j, \Sigma_i} p_j
  \end{split}
\end{equation}
Ricordando la definizione di $\alpha_j$ possiamo calcolare il commutatore:
\begin{equation}
  \begin{split}
    \sqbr{\alpha_j, \Sigma_i} &=
    \begin{pmatrix}
      0 & \sigma_j \\
      \sigma_j & 0
    \end{pmatrix}
    \begin{pmatrix}
      \sigma_i & 0 \\
      0 & \sigma_i
    \end{pmatrix} -
    \begin{pmatrix}
      \sigma_i & 0 \\
      0 & \sigma_i
    \end{pmatrix}
    \begin{pmatrix}
      0 & \sigma_j \\
      \sigma_j & 0
    \end{pmatrix} \\[8pt]
    &=
    \begin{pmatrix}
      0 & \sigma_j\sigma_i \\
      \sigma_j\sigma_i & 0
    \end{pmatrix} -
    \begin{pmatrix}
      0 & \sigma_i\sigma_j \\
      \sigma_i\sigma_j & 0
    \end{pmatrix} \\[8pt]
    &=
    \begin{pmatrix}
      0 & [\sigma_j, \sigma_i] \\
      [\sigma_j, \sigma_i] & 0
    \end{pmatrix} \\[8pt]
    &= 2i\epsilon_{jik}
    \begin{pmatrix}
      0 & \sigma_k \\
      \sigma_k & 0
    \end{pmatrix} = 2i\epsilon_{jik} \alpha_k
  \end{split}
\end{equation}
Ovvero:
\begin{equation}
  [H, S_i] = i \epsilon_{jik} \alpha_k p_j
\end{equation}
Otteniamo dunque che i singoli contributi di momento angolare e spin non sono conservati, ma possiamo definire il momento angolare totale come:
\begin{equation}
  \vec{J} = \vec{L} + \vec{S}
\end{equation}
Calcoliamo ora il commutatore tra l'hamiltoniana di Dirac e il momento angolare totale:
\begin{equation}
  \begin{split}
    [H, J_i] &= [H, L_i + S_i] = [H, L_i] + [H, S_i] \\[8pt]
    &= -i \epsilon_{ijl} \alpha_j p_l + i \epsilon_{jik} \alpha_k p_j \\[8pt]
    &= -i \epsilon_{ijl} \alpha_j p_l + i \epsilon_{ijl} \alpha_l p_j = 0
  \end{split}
\end{equation}
Dunque il momento angolare totale è una quantità conservata nell'equazione di Dirac.\\
Un'altra osservazione interessante riguarda le soluzioni dell'equazione di Dirac. Consideriamo il caso di una particella a riposo, ovvero con momento nullo $\vec{p} = 0$. Poichè il tensore quantità di moto è:
\begin{equation}
  p^\mu = \brackets{E, p_x, p_y, p_z} = \brackets{m, \pdv{}{x}, \pdv{}{y}, \pdv{}{z}} = i \brackets{\pdv{}{t}, \pdv{}{x}, \pdv{}{y}, \pdv{}{z}} = i \partial^\mu
\end{equation}
Otteniamo che l'equazione di Dirac diventa:
\begin{equation}
  \brackets{i\gamma^0 \partial_0 - m}\psi = 0 \implies \brackets{E \gamma^0 - m}\psi = 0
\end{equation}
Sostituendo la matrice $\gamma^0$ otteniamo:
\begin{equation}
  \begin{pmatrix}
    E & 0 & 0 & 0 \\
    0 & E & 0 & 0 \\
    0 & 0 & -E & 0 \\
    0 & 0 & 0 & -E
  \end{pmatrix}
  \begin{pmatrix}
    \phi_1 \\ \phi_2 \\ \phi_3 \\ \phi_4
  \end{pmatrix}
  = m \begin{pmatrix}
    \phi_1 \\ \phi_2 \\ \phi_3 \\ \phi_4
  \end{pmatrix}
\end{equation}
Ovvero stiamo cercando gli autovalori della matrice $E \gamma^0$: Troviamo dunque due autovalori positivi $E = m$ con autovettori:
\begin{equation}
    u_1 = N_1 \begin{pmatrix}
      1 \\ 0 \\ 0 \\ 0
    \end{pmatrix}, \qquad
    u_2 = N_2 \begin{pmatrix}
      0 \\ 1 \\ 0 \\ 0
    \end{pmatrix}
\end{equation}
E due autovalori negativi $E = -m$ con autovettori:
\begin{equation}
    u_3 = N_3 \begin{pmatrix}
      0 \\ 0 \\ 1 \\ 0
    \end{pmatrix}, \qquad
    u_4 = N_4 \begin{pmatrix}
      0 \\ 0 \\ 0 \\ 1
    \end{pmatrix}
\end{equation}
Ricordando che la funzione d'onda è del tipo:
\begin{equation}
  \psi = u \efunction^{-iEt}
\end{equation}
Otteniamo dunque le seguenti soluzioni per l'equazione di Dirac a riposo:
\begin{equation}
    \psi_1 = N_1 \begin{pmatrix}
      1 \\ 0 \\ 0 \\ 0
    \end{pmatrix} \efunction^{-imt}, \qquad
    \psi_2 = N_2 \begin{pmatrix}
      0 \\ 1 \\ 0 \\ 0
    \end{pmatrix} \efunction^{-imt}
\end{equation}
\begin{equation}
    \psi_3 = N_3 \begin{pmatrix}
      0 \\ 0 \\ 1 \\ 0
    \end{pmatrix} \efunction^{+imt}, \qquad
    \psi_4 = N_4 \begin{pmatrix}
      0 \\ 0 \\ 0 \\ 1
    \end{pmatrix} \efunction^{+imt}
\end{equation}
Le soluzioni con $E<0$ non possono essere semplicemente scartate in quanto non fisiche, in quanto l'equazione di Dirac è del primo ordine in tempo e dunque richiede 4 condizioni iniziali per essere risolta, una per ogni componente dello spinore.\\
Possiamo studiare l'equazione di dirac per sistemi non a riposo $\vec{p} \neq 0$. In questo caso otteniamo 4 soluzioni del tipo:
\begin{equation}
    u_1 = N_1\begin{pmatrix}
      1 \\[8pt] 0 \\[8pt] \dfrac{p_z}{E+m} \\[8pt] \dfrac{p_x + ip_y}{E+m}
    \end{pmatrix},
    u_2 = N_2\begin{pmatrix}
      0 \\[8pt] 1 \\[8pt] \dfrac{p_x - ip_y}{E+m} \\[8pt] -\dfrac{p_z}{E+m}
    \end{pmatrix},
    u_3 = N_3\begin{pmatrix}
      \dfrac{p_z}{E+m} \\[8pt] \dfrac{p_x + ip_y}{E+m} \\[8pt] 1 \\[8pt] 0
    \end{pmatrix},
    u_4 = N_4\begin{pmatrix}
      \dfrac{p_x - ip_y}{E+m} \\[8pt] -\dfrac{p_z}{E+m} \\[8pt] 0 \\[8pt] 1
    \end{pmatrix}
\end{equation}
Interpetazione di Dirac delle soluzioni negative: la particella può arrivare al più allo stato fondamentale, poiché gli stati ad energia negativa sono già occupati e per il principio di esclusione di Pauli la particella non può occupare quegli stati.\\
Questa interpretazione è detta \textbf{interpretazione del mare di Dirac}. In questo modello il passaggio da energia negativa a positiva viene mediante l'assorbimento di un fotone. Simmetricamente lo riempimento di una lacuna nel mare di elettroni provoca l'emissione di un fotone.\\
Dunque la produzione di una coppia elettrone-lacuna può essere vista come l'eccitazione di un elettrone dal mare di Dirac.\\
\subsection{Interpretazione con antiparticelle}
Nel 1933 C.D. Anderson scopre sperimentalmente l'esistenza del positrone, l'antiparticella dell'elettrone. Questa scoperta fornisce una conferma alle previsioni matematiche presentate da Dirac ed una valida alternativa all'interpretazione del mare di Dirac. Infatti possiamo interpretare le soluzioni negative come soluzioni per particelle con carica opposta, ovvero le antiparticelle.\\
Questo inoltre risolve il problema dell'incompatibilità dell'interpretazione del mare di Dirac con particelle bosoniche, che non rispettano il principio di esclusione di Pauli.\\
La scoperta sperimentale avvenne tramite l'osservazione di particelle cariche positivamente in una camera a nebbia immersa in un campo magnetico. Queste particelle presentavano la stessa massa dell'elettrone (le particelle più massive venivano bloccate da una lastra di piombo), ma percorso opposto a quello degli elettroni e, di conseguenza, carica opposta.\\
L'interpretazione con antiparticelle è detta di Feynman-Stueckelberg. Le antiparticelle possono essere interpretate come particelle che viaggiano indietro nel tempo. In questo modo l'annichilazione di una particella può essere descritta con il seguente diagramma di Feynman:

\addfeynman[Diagramma di Feynman per l'annichilazione di una particella e la sua antiparticella]{
    \draw[fermion]
        (-0.5,-0.5)node[below left]{$e^-$} -- (1,1) -- (3,1) -- (4.5,-0.5)node[below right]{$e^+$};
    \draw[photon]
        (1,1) -- (-0.5,2.5) node[above left]{$\gamma$};
    \draw[photon]
        (3,1) -- (4.5,2.5)node[above right]{$\gamma$};
    \draw[-stealth,thick] (-2,1-0.5) -- (-2,1+0.5) node[midway,left]{Time};
}{1.2}{feynman_annihilation}

Ovvero abbiamo la reazione:
\begin{equation}
    e^- + e^+ \rightarrow \gamma + \gamma
\end{equation}
Interpretando l'antiparticella come una particella che viaggia indietro nel tempo otteniamo:
\begin{equation}
  \psi = N \efunction^{-i\brackets{\vec{p}\cdot\vec{r} - Et}} \rightarrow N \efunction^{-i\vec{p}\cdot\vec{r}} \efunction^{i(-E)(-t)} = \psi
\end{equation}
Possiamo definire una quantità legata allo spin, valida per tutte le soluzioni e dunque non limitata a prendere $\hat{S}_z$ come riferimento.
\begin{definition}{Elicità}
    Definiamo l'operatore elicità come:
    \begin{equation}
        h \defineeq \dfrac{2\vec{S}\cdot\vec{p}}{\abs{\vec{p}}} = \dfrac{\boldsymbol{\Sigma}\cdot\vec{p}}{\abs{\vec{p}}}
    \end{equation}
\end{definition}
Questa quantità è detta \textbf{elicità} ed è una proprietà intrinseca delle particelle. Per particelle di massa nulla l'elicità è invariante per trasformazioni di Lorentz, mentre per particelle massive può essere cambiata da un sistema di riferimento ad un altro.\\
